\chapter{TTY layer, system calls and getty}
\label{cap:tty}

\section{From VGA renderer to virtual consoles}

Up to this point the VGA driver owned both the screen contents and the cursor
position. To support several Linux-like terminals, the code was split into two
layers:

\begin{itemize}
  \item \code{vga.rs} is a \emph{renderer}. It knows the 80$\times$25 cell
        matrix at \code{0xB8000}, the CRTC cursor registers and nothing about
        command lines or input;
  \item \code{tty.rs} owns eight virtual consoles, \code{tty0} through
        \code{tty7}. Each console has its own cell buffer, cursor and input
        queue.
\end{itemize}

Only the active console is copied to the hardware framebuffer. Switching
consoles is therefore a buffer copy plus a cursor update, exactly the same
idea used by the virtual terminal (\code{vt}) layer in Linux.

\begin{lstlisting}[caption={Per-console state in \code{tty.rs}.},label={lst:ttytty}]
#[derive(Clone, Copy)]
struct Tty {
    cells: [u16; vga::CELLS],
    row: usize,
    column: usize,
    input: [u8; INPUT_SIZE],
    head: usize,
    tail: usize,
}

static TTYS: SpinLock<[Tty; TTY_COUNT]> =
    SpinLock::new([Tty::new(); TTY_COUNT]);
static ACTIVE: AtomicUsize = AtomicUsize::new(0);
\end{lstlisting}

The input queue is a small circular buffer of 256 bytes. The \emph{producer}
is \code{pump\_input}, which drains the PS/2 keyboard and COM1; the
\emph{consumer} is the shell or a system call reading from the active console.

\section{Console switching}

The keyboard driver now exposes two kinds of results:

\begin{lstlisting}[caption={Keyboard result type.},label={lst:keyenum}]
pub enum Key {
    Char(u8),
    Console(usize),
}
\end{lstlisting}

The driver tracks Shift and Alt. When Alt is held, scancodes
\code{0x3B}..\code{0x42} (F1..F8) become \code{Key::Console(0..7)}. The TTY
layer translates that into \code{switch(index)} and redraws the selected
buffer. The same operation is available from the shell:

\begin{lstlisting}[style=shell,caption={Switching consoles from the shell.},label={lst:ttyshell}]
rinux> tty
console ativo: tty0 (disponiveis: tty0..tty7)
rinux> tty 3
rinux> tty
console ativo: tty3 (disponiveis: tty0..tty7)
\end{lstlisting}

Because the shell is a single synchronous loop, switching consoles while a
line is being edited discards the partial line and prints a fresh prompt on
the new console. This is a deliberate simplification: there is no
\code{task\_struct} per terminal yet.

\section{ANSI/VT100 escape sequences}

Output is not limited to printable characters. Each \code{Tty} keeps its own
small state machine, so a program can use the same escape sequences as a
Linux terminal without knowing which console is active:

\begin{lstlisting}[caption={Escape state stored in each console.},label={lst:ttyescape}]
enum EscapeState {
    Ground, // texto normal
    Start,  // recebeu ESC
    Csi,    // recebeu ESC [ e le parametros
}

struct Tty {
    // ...
    escape: EscapeState,
    csi_params: [u16; 4],
    csi_count: usize,
    csi_value: u16,
    csi_digits: bool,
}
\end{lstlisting}

\code{write\_byte} routes \code{0x1B} to the state machine. In the CSI state
the parser accepts decimal parameters separated by \code{;} and a final byte;
private markers such as \code{?} and intermediate bytes are consumed without
being printed:

\begin{lstlisting}[caption={Dispatch of ESC and CSI in \code{tty.rs}.},label={lst:ttywritebyte}]
fn write_byte(&mut self, byte: u8) {
    match self.escape {
        EscapeState::Ground => {
            if byte == 0x1B {
                self.escape = EscapeState::Start;
            } else {
                self.write_plain(byte);
            }
        }
        EscapeState::Start => {
            if byte == b'[' {
                self.escape = EscapeState::Csi;
                self.csi_params = [0; 4];
                self.csi_count = 0;
                // ...
            } else {
                self.escape = EscapeState::Ground;
            }
        }
        EscapeState::Csi => self.write_csi(byte),
    }
}
\end{lstlisting}

\begin{table}[H]
\centering
\caption{Supported CSI sequences.}
\label{tab:ansi}
\begin{tabularx}{\linewidth}{@{}lX@{}}
\toprule
\textbf{Sequence} & \textbf{Effect} \\
\midrule
\code{ESC[2J}, \code{ESC[3J} & Erases the whole display. \\
\code{ESC[J} & Erases from the cursor to the end of the display. \\
\code{ESC[K} & Erases from the cursor to the end of the line. \\
\code{ESC[H}, \code{ESC[<row>;<col>H} & Moves the cursor home or to a
1-based position, clamped to 80$\times$25. \\
\code{ESC[A}, \code{ESC[B}, \code{ESC[C}, \code{ESC[D} & Moves the cursor
up, down, right or left (default 1, or the given count). \\
\code{ESC[...m} & SGR is parsed and ignored; colors are not implemented yet. \\
\bottomrule
\end{tabularx}
\end{table}

The \code{clear} command is now ordinary output, not a special case in the
renderer:

\begin{lstlisting}[caption={Implementation of the \code{clear} command.},label={lst:ttyclear}]
pub fn clear() {
    write_bytes(b"\x1b[2J\x1b[H");
}
\end{lstlisting}

\code{ESC[2J} blanks the cells but, following the VT100 behavior, does not
move the cursor; the second sequence \code{ESC[H} returns it to the top-left
corner. Since \code{clear} goes through \code{write\_bytes}, the escape
sequence is also mirrored to COM1, so a serial terminal clears as well.
SGR colors, alternate screen, scrollback and cursor visibility are still out
of scope.

\section{System call ABI}

The gate \code{0x80} already existed for the ring 3 demos. It now dispatches
to \code{syscall.rs}, with the usual register convention:

\begin{table}[H]
\centering
\caption{Implemented system calls.}
\label{tab:syscalls}
\begin{tabularx}{\linewidth}{@{}clX@{}}
\toprule
\textbf{Number} & \textbf{Name} & \textbf{Arguments and result} \\
\midrule
0 & \code{exit} & \code{rdi} = status; does not return. \\
1 & \code{write} & \code{rdi} = buffer, \code{rsi} = length;
returns the number of bytes. \\
2 & \code{read} & \code{rdi} = buffer, \code{rsi} = length,
\code{rdx} = flags (bit 0 disables echo); returns the number of bytes read. \\
3 & \code{get\_tty} & no arguments; returns the active console index. \\
\bottomrule
\end{tabularx}
\end{table}

\subsection{User pointer validation}

A system call runs in ring 0 with the same page tables as the user. Touching an
unmapped user address would not generate a user page fault; it would be a
kernel access. Therefore \code{write} and \code{read} call
\code{paging::ensure\_user\_range} before dereferencing the buffer:

\begin{lstlisting}[caption={Validation and demand mapping of a user buffer.},label={lst:ensureuser}]
pub fn ensure_user_range(addr: u64, len: u64) -> Result<(), &'static str> {
    let end = addr.checked_add(len).ok_or("intervalo invalido")?;
    if addr < USER_BASE || end > USER_BASE + USER_REGION_SIZE {
        return Err("ponteiro fora da regiao do usuario");
    }
    // Mapeia com U/S = 1 cada pagina ausente de [addr, end).
    // ...
    Ok(())
}
\end{lstlisting}

The function rejects intervals outside the user region and materializes each
missing page with \code{PAGE\_USER | PAGE\_WRITE}. This is the didactic
equivalent of the careful \code{copy\_from\_user}/\code{copy\_to\_user}
helpers in Linux.

\subsection{Blocking read}

The \code{read} call is synchronous. While no byte is available it calls
\code{tty::pump\_input} and \code{spin\_loop}. The interrupt gate clears
\code{IF}, so the timer does not advance during the wait in the current
implementation. Printable bytes are stored in the user buffer and echoed when
the no-echo flag is clear; Backspace removes the last byte; Enter terminates
the line and returns the newline.

\section{getty and login in ring 3}

The first real user-mode program is a \emph{getty}/login. It is written in
assembly inside \code{user.rs}, in its own section, and copied as a
position-independent blob to \code{paging::USER\_CODE}:

\begin{lstlisting}[caption={Login flow of the ring 3 program.},label={lst:gettyasm}]
getty_start:
    GETTY_WRITE getty_login_msg, getty_login_end   # write(2)
    subq $128, %rsp                                # buffer na pilha
    movq %rsp, %rdi                                # buf
    movq $32, %rsi                                 # len
    xorq %rdx, %rdx                                # echo
    movq $2, %rax                                  # read
    int $0x80
    # compara "root" e pede "Password: "
    # segundo read com flags = 1 (sem eco)
    # compara "rinux" e escreve a mensagem final
    # exit(0) ou exit(1)
\end{lstlisting}

Using a dedicated section with code and strings keeps all RIP-relative
references valid after the copy. The read buffer lives on the user stack,
whose page is demand-mapped by the first system call that touches it.

The demonstration credentials are:

\begin{lstlisting}[style=shell,caption={Successful login.},label={lst:gettyout}]
rinux> getty
[GETTY] iniciando login em tty0...
rinux login: root
Password:
Bem-vindo ao rinux, root!
[GETTY] login concluido, voltando ao shell
\end{lstlisting}

\code{getty n} first switches to console \code{n} and then runs the program
there. In an automated serial test, the credentials must be sent only after
the login prompt of that console appears, because the serial port is a single
global stream.

\section{Relationship with Linux}

\begin{table}[H]
\centering
\caption{Correspondence with the Linux TTY subsystem.}
\label{tab:ttflinux}
\begin{tabularx}{\linewidth}{@{}lX@{}}
\toprule
\textbf{Linux} & \textbf{Here} \\
\midrule
\code{drivers/tty/tty\_io.c} & \code{src/tty.rs}: read/write queues. \\
\code{drivers/tty/vt/vt.c} & \code{src/tty.rs}: virtual consoles and
switching. \\
\code{drivers/video/console/vgacon.c} & \code{src/vga.rs}: cell renderer. \\
\code{drivers/tty/n\_tty.c} & Not implemented; lines are handled directly. \\
\code{arch/x86/entry/syscalls/syscall\_64.tbl} & \code{src/syscall.rs}. \\
\code{init/getty.c}, \code{login} & Raw ring 3 blob in \code{src/user.rs}. \\
\bottomrule
\end{tabularx}
\end{table}

\section{Current limitations}

\begin{itemize}
  \item There is one shell, not one shell per terminal. Switching consoles
        moves the same loop to another buffer.
  \item \code{getty} is not resident: it runs once, synchronously, and returns
        to the shell. A real system keeps a \code{getty} waiting on each TTY.
  \item The serial port is a single stream; keyboard input is what really
        distinguishes the consoles in interactive use.
  \item There is no line discipline: no canonical mode, signals, job control,
        \code{Ctrl+C}, \code{Ctrl+D} or terminal ioctls.
  \item \code{read} waits with interrupts disabled; the timer freezes during a
        login prompt.
\end{itemize}

\section{Kernel heap}

\code{src/heap.rs} adds the missing dynamic allocator. It places the heap at
\code{0xFFFFFFFFE0000000}, a separate PDPT from the kernel image and from the
demand-paging regions, and installs a \code{\#[global\_allocator]} so that
\code{alloc::vec::Vec}, \code{alloc::string::String} and the rest of the
\code{alloc} crate work inside the kernel.

The allocator is a first-fit free list with a 16-byte header immediately
before the returned pointer:

\begin{lstlisting}[caption={Block header and free-list node.},label={lst:heapblock}]
#[repr(C, align(16))]
struct Block {
    size: usize,
    next: *mut Block,
}
\end{lstlisting}

Blocks are kept in address order and coalesced when neighbours are free.
When the list cannot satisfy a request, the heap grows one 4 KiB page at a
time through \code{paging::alloc\_frame} and \code{paging::map\_page}. The
limit is 1 MiB (256 pages). There is no slab, buddy allocator, per-CPU cache
or NUMA support; the goal is to explain the interface and the growth path.
The \code{heap} applet allocates a \code{Vec<u64>} and a \code{String} and
prints the mapped/used/allocations counters.

\section{VFS and file descriptors}

\code{src/vfs.rs} introduces the two central objects of a Unix file layer:
the open file description (here a table entry with an offset and flags) and
the descriptor, an index into a fixed table of 16 entries. Descriptors 0, 1
and 2 are pre-opened on the console; \code{/dev/tty} and \code{/dev/console}
are accepted as console paths.

The flag and seek constants mirror the Linux ABI:

\begin{lstlisting}[caption={Subset of the VFS constants.},label={lst:vfsflags}]
pub const O_RDONLY: u64 = 0;
pub const O_WRONLY: u64 = 1;
pub const O_RDWR:   u64 = 2;
pub const O_CREAT:  u64 = 0x40;
pub const O_TRUNC:  u64 = 0x200;
pub const O_APPEND: u64 = 0x400;
pub const SEEK_SET: u64 = 0;
pub const SEEK_CUR: u64 = 1;
pub const SEEK_END: u64 = 2;
\end{lstlisting}

\code{open}, \code{close}, \code{read}, \code{write} and \code{lseek} are
implemented over the ramfs, the console and \code{/proc}. Reads update the
descriptor offset; \code{O\_APPEND} recomputes the end of the file on each
write; \code{lseek} supports \code{SEEK\_SET}, \code{SEEK\_CUR} and
\code{SEEK\_END}. There are still no inodes, dentries, mounts, permissions or
page cache, and the table is global rather than per process.

\section{The /proc pseudo-filesystem}

\code{src/proc.rs} implements a synthetic filesystem. Nothing is stored:
each read formats the current kernel state into a fixed 2 KiB buffer and
copies the requested slice according to the descriptor offset. The available
files are:

\begin{center}
\begin{tabular}{@{}ll@{}}
\toprule
\textbf{Path} & \textbf{Content} \\
\midrule
\code{/proc/version} & version, architecture and \code{\#1 SMP} \\
\code{/proc/uptime}  & seconds and centiseconds from the PIT ticks \\
\code{/proc/meminfo} & frames, heap and task counters \\
\code{/proc/ps}      & task table in the same format as the \code{ps} applet \\
\code{/proc/signals} & pending/blocked masks, state and exit status of tasks \\
\code{/proc/cpuinfo} & a small static CPU description \\
\code{/proc/mounts}  & ramfs, proc and devtmpfs lines \\
\code{/proc/cmdline} & boot command line \\
\bottomrule
\end{tabular}
\end{center}

\code{cat /proc/ps}, \code{ls -l /proc} and the \code{vfs} applet exercise the
same code path used by the ring 3 syscalls. There is no \code{seq\_file},
\code{proc\_dir\_entry}, permission check or namespace.

\section{Tasks and preemptive scheduler}

\code{src/task.rs} adds PID/PPID, task states and a \emph{weighted fair}
scheduler inspired by the Linux CFS. The shell is task 1 and runs on the boot
stack. \code{spawn(name, entry, nice)} creates a kernel thread with a 16 KiB
stack from the heap, a copy of the name, an entry function and a priority.
Each task has its own saved stack pointer and a start frame consumed by the
assembly routine below; besides the callee-saved registers, the routine saves
and restores \code{RFLAGS}, so the interrupt flag follows the task across a
switch:

\begin{lstlisting}[caption={Context switch with \code{RFLAGS}.},label={lst:ctxswitch}]
task_switch:
    pushfq
    pushq %rbx
    pushq %rbp
    pushq %r12
    pushq %r13
    pushq %r14
    pushq %r15
    movq %rsp, (%rdi)
    movq %rsi, %rsp
    popq %r15
    popq %r14
    popq %r13
    popq %r12
    popq %rbp
    popq %rbx
    popfq
    ret
\end{lstlisting}

A new task starts at \code{task\_trampoline}, which calls the entry function
and then \code{exit\_current}. \code{yield\_now} marks the current task ready
and lets the scheduler choose again, while \code{exit\_current} marks it
\code{Exited}. An exited task remains as a zombie until its parent calls
\code{wait}, which frees the slot; the next \code{spawn} then returns the old
stack to the heap.

The policy is weighted fair scheduling with \code{nice},
\code{sched\_prio\_to\_weight}-style weights and a virtual runtime. Every task
stores \code{nice} ($-20$ to $19$), its weight (from 88761 at $-20$ down to 15
at $19$, with 1024 at $0$) and \code{vruntime}. At each IRQ0 the current task
adds
\[
  \Delta \mathrm{vruntime} = \frac{1024 \ll 10}{\mathrm{weight}}
\]
to its virtual runtime, so a high-priority task accumulates virtual time more
slowly. \code{on\_tick} then chooses the runnable task (\code{Ready} or
\code{Running}) with the smallest \code{vruntime}; ties keep the current task
to avoid a needless switch. There is no fixed quantum: fairness emerges from
the per-tick accounting, and equal-weight tasks alternate roughly every tick
while a task with \code{nice} $-5$ receives about three times the CPU of a
\code{nice} $0$ task. New tasks and woken tasks are clamped to the minimum
\code{vruntime} of the runnable set, the same idea as \code{place\_entity} in
\code{kernel/sched/fair.c}, so they do not monopolize the CPU while catching
up. This is a didactic CFS: there is no red--black tree, \emph{lag} or EEVDF.

Because the timer handler itself can take the task-table lock, the spinlocks
of \code{sync.rs} follow the \code{irqsave} model: acquiring a lock disables
local interrupts and the guard restores the previous \code{RFLAGS} on drop.
Preemption is suspended while a ring 3 program is active, since that program
runs on the shared TSS interrupt stack instead of a per-task stack; the clock
and \code{vruntime} keep advancing. All tasks share the kernel address space;
there is no \code{fork}, \code{exec} or per-process page table. \code{wait}
exists as a parent/child operation over the same address space and reaps
zombies, as described below.

The scheduler also has a real blocking state. \code{task::sleep(ticks)} is the
equivalent of \code{schedule\_timeout}: the current task stores an absolute
deadline in \code{wake\_at}, changes to \code{Blocked} and switches away, so it
no longer competes for the CPU. At every IRQ0, \code{on\_tick} advances the
scheduler clock and moves expired tasks back to \code{Ready}; the idle loop
also executes \code{hlt} when no task is runnable, instead of spinning. The
timer wakeup path is the moral equivalent of the timer wheel used by
\code{schedule\_timeout}.

Blocking also supports events. \code{WaitQueue} mirrors the contract of
\code{wait\_queue\_head\_t} in \code{include/linux/wait.h}: a task records the
queue in \code{wait\_queue}, changes to \code{Blocked} and is only made
\code{Ready} when another context calls \code{wake\_up\_one} or
\code{wake\_up\_all}. \code{wait\_event} rechecks a condition after every
wakeup, and \code{wait\_event\_timeout} adds a deadline, like
\code{wait\_event\_timeout} in Linux. A woken task has its \code{vruntime}
raised to the current minimum before becoming runnable, so it cannot arrive
far behind the others. The queue identity is the address of the static
structure, and the waiters live in the task table, which is a simplified
version of Linux's intrusive linked lists.

The \code{spawn [nome] [n] [cpu|sleep|wait|timeout|signal|exit] [nice]} applet
creates up to eight workers. In the default \code{cpu} mode each worker yields
once and then consumes 10 real CPU ticks per round, measured by the per-task
\code{cpu\_ticks} counter, so the observed interleaving comes from preemption
and the same work takes longer at a lower priority. With \code{sleep}
(\code{spawn nap 2 sleep -5}) the worker calls \code{task::sleep} between
rounds. With \code{wait}, it sleeps on the event queue until the
\code{notify} applet calls \code{wake\_up\_all} (or \code{notify one} for
\code{wake\_up}); \code{timeout} uses \code{wait\_event\_timeout} and reports
whether the event arrived within one second; \code{signal} installs a
\code{SIGUSR1} handler and sleeps; \code{exit} ends with status 42 after half
a second, leaving a zombie for \code{wait}. In all blocking modes \code{ps}
shows the task as \code{dormindo} (or \code{parada} after \code{SIGSTOP})
while the shell stays responsive. The \code{nice [pid] [valor]} applet reads
or changes a task's priority, like \code{setpriority}. \code{ps} prints PID,
PPID, state, \code{NICE}, \code{CPU} ticks, \code{VRUNTIME}, yield count,
preemption count (\code{PREEMP}) and name. Two equal workers finish with the
same CPU ticks and \code{vruntime}; a worker created with \code{nice} $-5$
finishes well before an equal \code{nice} $0$ worker and accumulates much
less \code{vruntime} per tick.

\section{Signals and process lifecycle}

Signals follow the Unix model in a compact form. Each task stores a 32-bit
\code{pending} mask, a \code{blocked} mask and an array of 32
\code{SignalAction} values, numbered from 1 (traditional Unix signals). The
disposition is one of:

\begin{lstlisting}[caption={Signal dispositions.},label={lst:signals}]
pub enum SignalAction {
    Default,
    Ignore,
    Handler(fn(i32)),
}
\end{lstlisting}

\code{Default} applies the classic action: terminate, stop or ignore,
depending on the signal. \code{SIGKILL} (9) and \code{SIGSTOP} (19) are
special: they cannot be blocked, ignored or have their disposition changed.
\code{SIGCONT} (18) always removes a task from \code{Stopped}, even when the
signal itself is blocked. \code{SIGCHLD} (17) and \code{SIGURG}/\code{SIGWINCH}
are ignored by default; \code{SIGTSTP}/\code{SIGTTIN}/\code{SIGTTOU} stop by
default; the remaining signals terminate by default.

Delivery happens only at safe points, never inside an IRQ handler, because a
handler may sleep or switch tasks. \code{check\_signals} is called from
\code{yield\_now}, \code{sleep}, \code{block\_current}, \code{wait\_child}, the
shell loop/idle path and the CPU worker loop. It takes the lowest pending
signal not in the blocked mask, clears its bit and executes the disposition.
The \code{in\_handler} flag prevents recursive delivery while a handler runs;
a signal that arrives during a handler waits for the next safe point.

Sending a signal goes through \code{task::kill(pid, sig)}. For the current
task the bit is only queued. For another task, \code{SIGKILL}, \code{SIGSTOP}
and default-terminating signals are applied immediately, since the target is
not executing and there is no kernel context to save; handler signals are
queued and a blocked target is woken so it reaches a safe point. A signal
handler is a kernel function pointer, not a user address: it runs in the
context of the target task, with the same kernel stack and address space.
\code{kill(pid, 0)} only checks existence and therefore succeeds for a
zombie.

The applets mirror the API: \code{kill} sends a number or a name;
\code{sigaction} reads or installs the disposition; \code{sigblock} changes
the mask; \code{sig} prints pending/blocked masks plus the exit status; and
\code{/proc/signals} exposes the same table. Handler signals are queued and
delivered at a safe point; default terminating or stopping signals sent to a
non-current task take effect immediately.

Lifecycle is tied to the same table. A task that calls
\code{exit\_current\_with\_status(status)} becomes \code{Exited}, stores its
status and, when killed by a signal, the fatal signal number
(\code{128 + sig}). Its parent receives \code{SIGCHLD} and, if it was blocked
in \code{wait}, is woken. \code{wait\_child(target)} waits on a global
\code{CHILD\_QUEUE}; when a zombie child is present it returns a
\code{WaitResult} (pid, status, signal) and marks the slot \code{Free}. A
ten-tick timeout covers the race between checking the condition and sleeping,
so a wakeup cannot be lost. Until reaped, a zombie keeps its PID and can
fill the 16-slot table; calling \code{wait} is therefore required before
respawning many short-lived children.

\section{New system calls}

The \code{int 0x80} table was extended with file and process calls:

\begin{center}
\begin{tabular}{@{}cll@{}}
\toprule
\textbf{Number} & \textbf{Name} & \textbf{Arguments} \\
\midrule
0  & \code{exit}     & \code{status} \\
1  & \code{write}    & \code{buf, len} \\
2  & \code{read}     & \code{buf, len, flags} \\
3  & \code{get\_tty} & -- \\
4  & \code{open}     & \code{path, len, flags} \\
5  & \code{close}    & \code{fd} \\
6  & \code{read\_fd} & \code{fd, buf, len} \\
7  & \code{write\_fd}& \code{fd, buf, len} \\
8  & \code{lseek}    & \code{fd, offset, whence} \\
9  & \code{getpid}   & -- \\
10 & \code{getppid}  & -- \\
11 & \code{kill}     & \code{pid, sig} \\
\bottomrule
\end{tabular}
\end{center}

\code{filetest} runs a ring 3 blob that opens \code{/etc/motd} with
\code{open}, reads it with \code{read\_fd}, writes the content with
\code{write\_fd} and closes the descriptor. \code{getpid} and \code{getppid}
read the current task from the scheduler instead of returning constants.
\code{kill} validates the signal number and delegates to \code{task::kill};
\code{sig = 0} only checks that the target exists.
